\section{Experimental Results (40 points, 600 words)}
\label{sec:experimentalresults}
% \subsection{Datasets}
% Include a description of the data used, what it contains, train test splits, etc.
% \subsection{Implementation Details}
% Include a description of the design choices, hyperparameters, and other implementation specific details.

% \subsection{Results}
% Analyze and discuss your results. Table \ref{tab:example-table} shows an example of a table.
% Include plots and visualizations such as Figure \ref{fig:one-col-figure}.

% % \subsection{Discussion}
% % Additionally, you can include a separate discussion section in which you summarize the results, indicate possible limitations of your work, and provide suggestions for future research.

% \subsection{Citing references}
% LaTeX manages the citations for you. You simply have to add the bibtex entry for the reference in the references.bib file and you can use the citation in the paper \cite{langley00}. You can also cite multiple references in one command \cite{DudaHart2nd,Newell81,kearns89}. The bibtex entries for a paper can be obtained in websites such as google scholar or dblp.

% \begin{table}[ht]
%     \centering
%     \label{tab:example-table}
%     \caption{Example Table}
%     \begin{tabular}{lllll}
%         \toprule
%         \multirow{2}{*}{Models} & \multicolumn{3}{c}{Metric 1} & Metric 2\\
%         \cmidrule{2-4} \cmidrule{5-5} \\
%         {} & precision & recall & F-score  & R@10 \\
%         \midrule
%         model 1 & 0.67  & 0.8 & 0.729  & 0.75 \\
%         model 2 & 0.8 & 0.9 & 0.847 & 0.85 \\
%         \bottomrule
%     \end{tabular}
% \end{table}

% \textbf{Points:} \textit{Experiments, analysis, and evaluation of the methods on data, benchmarks, simulations}
% \begin{itemize}
%     \item Does the evaluation contain proper baselines?
%     \item Is the evaluation correctly performed from a data analytics/machine learning perspective (depending on the field)
%     \item Are there ablation experiments?
%     \item Error analysis (e.g. analysis on which cases does it work and which cases does it fail?, is there a pattern across the errors? Analysis on the root cause of the problems?, etc.)
%     \item Other experiments (e.g. interpretability? Attention visualization? Qualitative?, etc.)
%     \item Were there statistical analysis?
% \end{itemize}


% = PARAGRAPH 1 - Baseline comparison

Present the aggregate BLEU scores from your run (BLEU-1: 0.468, BLEU-2: 0.286, BLEU-3: 0.189, BLEU-4: 0.130) and compare them directly with Table 7 in \citet{hwang2021cometatomic}. This serves as your baseline sanity check \textemdash{} confirming your evaluation pipeline reproduces the paper's numbers before diving into the per-relation breakdown. Note that your overall scores match, which validates the experimental setup

% = PARAGRAPH 2 - Category-level results

Present the three-way category table and not that differences are suprisingly small (BLEU-1 ranging only from 0.708-0.763). Discuss the ordering: social-iteractions leads on BLEU-1 and BLEU-2, while event-centred leads on BLEU-3 and BLEU-4. This is where you address the original hypothesis and not it was only partially supported.

% = PARAGRAPH 3 - Per-relation results: strong performers

Focus on \verb|HinderedBy| (0.891) and \verb|ObjectUse| (0.876) as standout relations. For \verb|HinderedBy| specifically, connetc back to the paper's finding that zero-shot GPT-2 scored only 1.3\% plausibility on this relation \textemdash{} this contrast is your strongest qualitative insight and directly addresses the error analysis requirement.

% = PARAGRAPH 4 - Per-relation results: weak performers

Discuss \verb|Causes| (0.315), \verb|HasProperty| (0.428), and \verb|CapableOf| (0.462) as the weakest relations. Offer a possibly explanation \textemdash{} these tend to require more open-ended, diverse outputs, making exact n-gram overlap with reference tails inherently harder to achieve. This is your error analysis paragraph.

% = PARAGRAPH 5 - Statistical caveat / ablation note

Flag the class inbalance issue explicitly: the weak performers also have very low counts (\verb|Causes|: 9, \verb|Desires|: 1, \verb|NotDesires|: 1), meaning their scores are unreliable estimates. Mention that \verb|xReason| and \verb|NotDesires| each have only 1 test example and should be disregarded. This addresses the statistical analysis requirement from the rubric and is an honest limitation of the per-relation breakdown.

% = PARAGRAPH 6 - Why category-level aggregation masks per-relation variance

Tie paragraph 4 and 5 together by explaining that because weak relations have very low counts, they contribute little weight to the category averages \textemdash{} which is why category-level BLEU scores look deceptively uniform. This is your ablation-style insight: the coarse-grained view hides what the fine-grained view reveals, which is the core argument of your whole investigation (my approach is the coarse grained one btw, apparently it sucks)
